\section{Maximal roofwise capacity and the microscopic residual}
\label{sec:v108-optimality}
The preceding pool has a precise maximality property. We specify the
class of certificates, since no optimality over all conceivable
physical transports is asserted.

\begin{definition}[A roofwise window certificate]
\label{def:v108-certificate}
Fix the same $r,\Delta,K$ and the same admissible comparison space
as in Definition~\ref{def:v108-pool}. A certificate for an accepted
source height $Z(t)$ consists of nonnegative measurable weights
$w_+(t,\theta),w_-(t,\theta)$ with $w_++w_-\le1$.
The positive weight is supported on segments of length less than
$\Delta$ containing $[t,t+r]$ in their interior, and the negative
weight on those containing $[t-r,t]$. Require
\begin{equation}\label{eq:v108-certificate}
 0\le Z(t)\le U(t),\qquad
 Z(t)\le K\int (w_++w_-)p_\vee\,d\theta.
\end{equation}
Weights may depend measurably on the original roof. No selector is
transported, no derivative of a selector is used, and the comparison
side has the same radius, collision count and exact labels.
\end{definition}

\begin{theorem}[Attained maximality]\label{thm:v108-maximal}
Among all certificates in Definition~\ref{def:v108-certificate},
the largest accepted height, pointwise almost everywhere, is
$V_{r,K}$. It has an explicit measurable certificate. Thus
$R_{r,K}$ is the smallest residual in this specified class.
\end{theorem}
\begin{proof}
Admissible supports and the one-copy constraint imply
$w_++w_-\le Q_r$. Therefore \eqref{eq:v108-certificate} gives
$Z\le\min\{U,KP_r\}=V_{r,K}$.
For attainment, on $P_r>0$ set
$\gamma(t)=V_{r,K}(t)/(KP_r(t))$, and set $\gamma=0$ on $P_r=0$.
Use $w_+=\gamma q_+$ and $w_-=\gamma q_-$.
These are measurable, their sum is at most one, and their certified
capacity is exactly $V_{r,K}$. When $P_r=0$ both sides vanish.
The source allocation is the distinct multiplier $\beta_{r,K}$
in \eqref{eq:v108-source-split}. This proves attainment without a
selection theorem or an optimization over an uncountable list of
angular families.
\end{proof}

\begin{theorem}[Domination of fixed-family allocations]
\label{thm:v108-dominance}
Let a revision-107 family act on the same source $U$, with comparison
lengths below $\Delta$, window width at least $r$, overlap at most
$N$, and block budget $\kappa$. Its accepted source $V^{\rm old}$
then satisfies
\begin{equation}\label{eq:v108-dominance}
 V^{\rm old}(t)\le V_{r,N\kappa}(t),\qquad
 R_{r,N\kappa}(t)\le U(t)-V^{\rm old}(t).
\end{equation}
The statement holds for every such family; no family maximizing the
old blockwise criterion has to exist or be selected.
\end{theorem}
\begin{proof}
At a roof $t$, retain only blocks with $T_j(t)>0$.
Lemma~\ref{lem:v107-window} puts their comparison supports inside
$\{Q_r(t,\cdot)=1\}$. The blockwise fraction obeys
$\beta_jT_j(t)\le\beta_jC_j\le\kappa P_j$.
For a finite collection of active blocks,
\[
 \sum_j\beta_jT_j(t)\le\kappa\int
       p_\vee(\theta)\sum_j w_j(\theta)\,d\theta
       \le N\kappa P_r(t).
\]
The full inequality follows by monotone convergence and the inherited
common versions for the countable family. Since source sets are
disjoint and $0\le\beta_j\le1$, also $V^{\rm old}\le U$.
Taking the minimum proves the first inequality; subtracting from
$U$ proves the second. This is domination at the original roof,
not merely an inequality between total angular capacities.
\end{proof}

\begin{proposition}[The new certificate is strictly broader]
\label{prop:v108-strict}
There are regular interval-source data for which no nonzero
whole-source-window block allocation of width $r$ is possible but
the roofwise certificate accepts the entire selected source on a
roof interval of positive length.
\end{proposition}
\begin{proof}
Take $r=1/10$, $s=1/1000$, and two disjoint angular sets of measure
one, each with constant minimum and maximum trace equal to one.
On the first set the physical interval is $(0,s)$; on the second it
is $(3s/8,r+7s/8)$. Select only the first set and take
$\Delta>r+s/2$, $K=1$. For $3s/8<t<7s/8$ the second interval
contains $[t,t+r]$ in its interior. Hence $U(t)=1$, $P_r(t)\ge1$
and $V_{r,1}(t)=1$ on that roof interval.
Neither available interval has length at least $r+s$. Consequently
neither can contain a translate by $[0,r]$ or $[-r,0]$ of the
entire first interval. Subdividing its angular set does not alter
its interval. Thus no nonzero revision-107 block comparison is
available for this selected source at width $r$.
All traces are constant, so the endpoint variation is zero.
These data demonstrate strictness of the certificate classes. They
are not asserted to constitute a new Lorentz realization or an
asymptotic example for the original collision process.
\end{proof}

Return now to the original unmatched source
$U=\mathsf U_B^{\rm micro}$ of Theorem~\ref{thm:v105-original}.
Its selector $\widetilde H_B$ and its restriction $0<\ell<h_B$
are unchanged. Set
\begin{equation}\label{eq:v108-schedule}
 \Delta_B=8\delta(B),\qquad
 0<r_B<\Delta_B,\qquad
 1\le K_B\le C\varepsilon(B)^{-1/32}.
\end{equation}
Here $r_B$, $h_B$ and $K_B$ are fixed at each outer stage before
the collision limit, uniformly over the radius and exact labels.
Define the canonical residual on the original source by
\begin{align}
 \mathsf R_B^{\rm can}(t)
   &=\bigl(\mathsf U_B^{\rm micro}(t)-K_BP_{r_B}(t)\bigr)_+
                                                    \label{eq:v108-residual}\\
   &=\int c^{\varepsilon(B)}(\theta)\one_{\{a(\theta)<t<b(\theta)\}}
        \widetilde H_B(\Psi(\theta,t))
                   (1-\beta_{r_B,K_B}(t))\,d\theta.\notag
\end{align}
The residual does not depend on a choice of angular source blocks.
It retains every original positive weight not paid by the complete
admissible roofwise capacity pool.

\begin{theorem}[Canonical reduction of the unchanged assertion]
\label{thm:v108-original}
Under the inherited scalar and physical-source inputs, and the
older budgets of Theorem~\ref{thm:v105-original}, the construction
\eqref{eq:v108-schedule} satisfies
\begin{equation}\label{eq:v108-original-error}
 \mathcal E_M\le\mathcal H_M(\mathsf R_B^{\rm can})
                                          +C_MB^{-1/384}.
\end{equation}
The original assertion \eqref{eq:v76-target} is equivalent to
\begin{equation}\label{eq:v108-final-height}
 \mathcal H_M(\mathsf R_B^{\rm can})\longrightarrow0
                         \quad\hbox{for every fixed }M.
\end{equation}
At the same or a smaller fixed window and the same total cost,
$\mathsf R_B^{\rm can}$ is pointwise no larger than any
revision-107 residual covered by Theorem~\ref{thm:v108-dominance}.
\end{theorem}
\begin{proof}
Corollary~\ref{cor:v108-allocation} decomposes the complete
finite-count microscopic source into its accepted and residual
parts. The inherited schedule has $\delta=O(\varepsilon^4)$.
Therefore Theorem~\ref{thm:v108-pool} gives
\[
 \mathcal H_M(V_{r_B,K_B})\le
 C\varepsilon^{-1/32}\{
            \varepsilon^{1/16}+O(\varepsilon^3)\}
       =O(\varepsilon^{1/32})=O(B^{-1/384}).
\]
The inherited scalar reduction and subadditivity of
$\mathcal H_M$ prove \eqref{eq:v108-original-error}.
The displayed decay implies the original assertion by sending
$B$ to infinity. Conversely the original assertion implies
$\mathcal H_M(\mathsf U_B^{\rm micro})\to0$ by
Theorem~\ref{thm:v105-original}; use
$0\le\mathsf R_B^{\rm can}\le\mathsf U_B^{\rm micro}$.
The pointwise comparison is Theorem~\ref{thm:v108-dominance}.
\end{proof}

\begin{proposition}[Finite-count exhaustion and its quantifiers]
\label{prop:v108-exhaustion}
Fix the collision count, radius, exact labels and $\Delta$.
Suppose the selected source is supported on the admissible
comparison space with $0<\ell<\Delta$, and $K\ge1$.
Then $R_{r,K}(t)$ decreases to zero for almost every $t$ as
$r\downarrow0$. This conclusion alone does not establish
\eqref{eq:v108-final-height}.
\end{proposition}
\begin{proof}
For every interior point of a positive-length segment, a sufficiently
small one-sided window lies in that segment. Hence, as $r\downarrow0$,
$Q_r(t,\theta)$ increases to
$\one_{\{0<\ell(\theta)<\Delta,\ a(\theta)<t<b(\theta)\}}$.
Monotone convergence and $p_\vee\ge c^\varepsilon$ give
$\lim_{r\downarrow0}P_r(t)\ge U(t)$.
The positive-part formula proves the asserted finite-count limit.
Here the collision count was fixed first. In
\eqref{eq:v108-final-height}, $r_B$ must be fixed first, followed by
the complete physical sum, the prescribed essential suprema and
collision limsup, with $B$ outermost. No interchange of these two
orders has been proved or used.
\end{proof}

The canonical construction removes the need to choose angular
blocks and to prove a separate overlap bound for their comparison
supports. It pays the complete endpoint trace loss and exhausts the
available capacity within the stated window-certificate class.
It does not prove that this capacity has a collision-uniform surplus
on the original microscopic source. In particular
\eqref{eq:v108-final-height} remains a substantive estimate on that
source, not a consequence of finite-count exhaustion. Nor does the
construction verify additional singular trace extensions outside
the inherited physical-source inputs. The original raw pointwise
assertion, rather than a reformulated comparison problem, remains
the organizing assertion of the paper.
